\chapter{What Is Doxonomics?}
\label{ch:what-is-doxonomics}

\epigraph{The most potent weapon in the hands of the oppressor is the mind of the oppressed.}{Steve Biko}

\noindent
Human beings do not merely inhabit economies and legal systems.
They inhabit interpretive worlds: shared assumptions about what people are, what is possible, what is deserved, and what cannot be changed.
These assumptions are rarely spontaneous.
They are funded, repeated, institutionalized, and defended.

\textbf{Doxonomics}\index{doxonomics} is the study of how material power shapes public belief.
The term derives from the Greek \textit{doxa} (opinion, common belief, public orthodoxy) and denotes the systematic investigation of how beliefs become socially normal.
Its central question is not whether a belief is true, but:

\begin{quote}
\textit{What did it cost to make this belief appear self-evident?}
\end{quote}

This question is not rhetorical.
It is a research question with empirical content.
Global advertising expenditure surpassed one trillion dollars in 2024.
U.S.\ lobbying and political spending exceed tens of billions annually.
Philanthropic grants to policy-oriented think tanks run into hundreds of millions per year, channeled through networks of foundations whose 990 filings are public record.
University curricula shape the foundational assumptions of millions of students per decade.
Platform algorithms determine which narratives reach which audiences at what frequency.
Each of these is a measurable flow of resources into the production of belief, and none of them is systematically tracked by any existing discipline as a \emph{unified} phenomenon.

Doxonomics proposes to track them.


\section{The Object of Study}
\label{sec:ch1-object}

Modern societies devote enormous resources to the production of goods, laws, and infrastructures.
They also devote immense resources to the production of legitimacy, resignation, aspiration, and moral common sense.
A society does not only distribute wealth; it distributes interpretations of the world.
Some interpretations are amplified until they appear obvious.
Others are starved of attention until they appear unrealistic or unintelligible.

Doxonomics studies this process.
It asks how money becomes belief, how institutions become credibility, how repetition becomes reality, and how narratives become embedded in everyday life \autocite{bourdieu1991, herman1988}.

Consider a concrete example.
The proposition ``human nature is selfish'' is not merely a philosophical claim.
It is an anthropological assumption embedded in economics curricula (taught to over ten million students per decade in the United States alone), management theory (principal-agent models presuppose worker self-interest), policy design (incentive-based regulation assumes that agents will not cooperate without material reward), advertising (which addresses the consumer as an isolated desire-maximizer), and popular culture (which narrates life as competition).
Each of these institutional sites has a budget.
Each produces narrative output.
Each reaches an audience.
The combined effect is an interpretive environment in which ``selfishness is natural'' operates as background reality rather than as a contestable claim.

Is this belief \emph{true}?
That is an empirical question, and the experimental evidence is considerably more mixed than the institutional consensus suggests (Chapter~\ref{ch:human-nature}).
The doxonomic question is different: \emph{what institutional apparatus sustains this belief, at what cost, through which channels, and with what social effects?}

\begin{definition}[Doxonomic System]
\label{def:doxonomic-system}
\index{doxonomic system}
A \emph{doxonomic system} is a tuple
\[
  \doxsys = (\agentspace, \insteco, \narr, \beliefspace, \Phi, \Psi)
\]
where:
\begin{itemize}[nosep]
  \item $\agentspace$ is a set of agents (funders, intermediaries, audiences),
  \item $\insteco$ is an institutional ecology (the ensemble of belief-processing institutions),
  \item $\narr$ is a narrative space (the set of available narratives and framings),
  \item $\beliefspace$ is a belief space (the space of possible population-level belief configurations),
  \item $\Phi: \R^m_{\geq 0} \times \insteco \to \narr$ is the \emph{production function} mapping funding and institutions to narrative outputs,
  \item $\Psi: \narr \times \beliefspace \to \beliefspace$ is the \emph{internalization function} mapping narrative exposure and prior beliefs to updated beliefs.
\end{itemize}
\end{definition}

\begin{remark}
This definition is deliberately abstract.
It does not specify the internal structure of agents, the topology of the institutional ecology, or the dimensionality of the belief space.
These are specified in later chapters as needed.
The purpose of the tuple is to establish that doxonomics has a \emph{formal object}: a mathematical structure that can be analyzed, compared, and measured.
\end{remark}

The aim is not conspiracy thinking.
A doxonomic system need not have a single controller.
Coordination may be deliberate, distributed, emergent, or structural.
A belief can dominate because many institutions, each acting locally and rationally, produce the same directional effect.
\textcite{herman1988} demonstrated this for media: five structural filters (ownership, advertising dependence, sourcing, flak, and ideological convergence) produce systematic bias without requiring any explicit coordination.
Doxonomics generalizes this insight beyond media to the full institutional landscape of belief production.

In this respect, doxonomics is closer to epidemiology than to melodrama.
Epidemiology does not require a single villain who spreads disease.
It maps transmission pathways, identifies risk factors, measures incidence, and designs interventions.
Doxonomics does the same for belief.


\section{The Core Conversion Problem}
\label{sec:ch1-conversion}

The foundational problem of doxonomics is the \textbf{conversion problem}\index{conversion problem}:

\begin{quote}
\textit{How are resources converted into durable public belief?}
\end{quote}

This conversion is never one-to-one.
One million dollars spent on direct advertising is not equivalent to one million dollars spent on textbooks, academic chairs, or platform placement.
Different channels have different reach, credibility, persistence, and downstream effects.
A think-tank report that shapes a government white paper that shapes a curriculum guideline that shapes twenty years of classroom instruction has a different doxonomic trajectory than a viral social media post that is forgotten in a week.

\begin{model}[First-Order Belief Equation]
\label{mod:first-order}
\index{belief equation}
As a first approximation, the social strength of belief $i$ at time $t$ is
\begin{equation}
\label{eq:first-order}
  \belief{i}(t) = \sum_{k=1}^{m} \funding{k}(t) \cdot \credibility{k} \cdot \relevance{k}{i} \cdot \reach{k} \cdot \persistence{k} - \counter{i}(t)
\end{equation}
where $\funding{k}(t)$ is funding through channel $k$, $\credibility{k}$ is the credibility weight of that channel, $\relevance{k}{i}$ is the relevance of channel $k$ to belief $i$, $\reach{k}$ is the population fraction reached, $\persistence{k}$ is the temporal persistence factor, and $\counter{i}(t)$ is the aggregate countervailing pressure against belief $i$.
\end{model}

This equation is a conceptual scaffold, not a mature model.
It is linear, static, and treats parameters as exogenous constants when in reality they are endogenous, time-varying, and interdependent.
Its purpose is to formalize a basic claim: beliefs are not just argued; they are \emph{financed exposures mediated by institutions}.
Later chapters will replace this linear approximation with dynamical models (Chapter~\ref{ch:formal-models}), threshold effects (Chapter~\ref{ch:common-sense-capture}), and network diffusion (Chapter~\ref{ch:network-analysis}).

Even in its crude form, the equation reveals something important.
Consider two beliefs:
\begin{itemize}[nosep]
  \item $p_1$: ``human nature is selfish,'' promoted through economics curricula ($\credibility{} \approx 0.9$, $\persistence{} \approx 30$ years), management training, advertising, and entertainment.
  \item $p_2$: ``humans are conditional cooperators,'' supported by experimental evidence but with minimal institutional infrastructure ($\funding{} \approx 0$, $\persistence{} \approx 0$).
\end{itemize}
The first-order model predicts $\belief{1} \gg \belief{2}$ regardless of the evidential merits of $p_1$ versus $p_2$.
This is the asymmetry that doxonomics studies.

\begin{remark}
The equation should not be read as claiming that \emph{only} money matters.
Evidence, lived experience, peer discussion, and individual reasoning all contribute to belief formation.
But the equation identifies a structural advantage: beliefs with institutional support have a baseline ``subsidy'' that beliefs without institutional support lack.
The playing field of ideas is not level; doxonomics measures the tilt.
\end{remark}


\section{The Belief Production Chain}
\label{sec:ch1-chain}

The full process from expenditure to social consequence is the \textbf{belief production chain}\index{belief production chain}.
It consists of seven stages, each with its own logic, costs, and failure modes:

\begin{enumerate}[nosep]
  \item \textbf{Funding}: resources are allocated to belief-production channels.
  \item \textbf{Institutional processing}: institutions (universities, think tanks, media, PR firms) transform funding into credentialed narrative.
  \item \textbf{Narrative output}: the processed material takes the form of research papers, op-eds, textbooks, films, curricula, advertisements, policy briefs.
  \item \textbf{Circulation}: narratives reach audiences through media, education, social networks, and platform algorithms.
  \item \textbf{Internalization}: audiences update their beliefs, norms, and practical dispositions in response to narrative exposure.
  \item \textbf{Behavioral consequence}: internalized beliefs alter behavior: cooperation, voting, consumption, labor compliance, tolerance of inequality.
  \item \textbf{Reproduction}: behavioral changes generate new institutional demands, electoral outcomes, and philanthropic priorities that feed back into the funding stage.
\end{enumerate}

\begin{center}
\begin{tikzpicture}[
  node distance=1.8cm,
  every node/.style={draw, rounded corners, minimum height=1cm, minimum width=2.2cm, align=center, font=\small},
  arrow/.style={-Stealth, thick}
]
  \node (F) {Funding};
  \node (I) [right=of F] {Institutional\\Processing};
  \node (N) [right=of I] {Narrative\\Output};
  \node (C) [right=of N] {Circulation};
  \node (Int) [below=1.2cm of C] {Internalization};
  \node (B) [left=of Int] {Behavioral\\Consequence};
  \node (R) [left=of B] {Reproduction};

  \draw[arrow] (F) -- (I);
  \draw[arrow] (I) -- (N);
  \draw[arrow] (N) -- (C);
  \draw[arrow] (C) -- (Int);
  \draw[arrow] (Int) -- (B);
  \draw[arrow] (B) -- (R);
  \draw[arrow, dashed] (R) -- (F) node[midway, draw=none, minimum height=0pt, minimum width=0pt, above, font=\footnotesize\itshape] {feedback};
\end{tikzpicture}
\end{center}

The dashed feedback arrow is critical.
The chain closes when the beliefs produced by the system generate exactly the behaviors and institutions that sustain those beliefs.
At that point the system has reached a \emph{doxonomic equilibrium}.
At equilibrium, the belief appears to reproduce itself spontaneously, without visible effort.
This is the condition that Bourdieu called \textit{doxa}: a belief so deeply embedded that it is no longer experienced as a belief at all.

Chapter~\ref{ch:belief-production-chain} formalizes each stage as a mathematical operator and the full chain as their composition.

\begin{example}[The Neoliberal Belief Production Chain]
\label{ex:neoliberal-chain}
The belief ``markets are the most efficient allocators of resources'' can be traced through a concrete chain:
\begin{enumerate}[nosep]
  \item \textbf{Funding}: Koch, Scaife, Olin, and Bradley foundations directed hundreds of millions of dollars into free-market intellectual infrastructure from the 1970s onward \autocite{mayer2016}.
  \item \textbf{Institutional processing}: Heritage Foundation, Cato Institute, American Enterprise Institute, and Mont P\`{e}lerin Society produced research, policy briefs, and expert testimony \autocite{mirowski2009}.
  \item \textbf{Narrative output}: op-eds, books (\textit{Free to Choose}, \textit{The Road to Serfdom}), textbooks (Mankiw's \textit{Principles of Economics}), and television programs.
  \item \textbf{Circulation}: editorial pages of \textit{Wall Street Journal}, \textit{Financial Times}; economics curricula at thousands of universities; cable news; social media.
  \item \textbf{Internalization}: generations of students, policymakers, journalists, and voters absorbed market-efficiency assumptions as common sense.
  \item \textbf{Behavioral consequence}: deregulation, privatization, austerity, reduced union membership, increased tolerance for inequality.
  \item \textbf{Reproduction}: beneficiaries of deregulation and low taxation reinvested in the same think tanks and foundations, closing the loop.
\end{enumerate}
This is not a conspiracy theory.
It is a documented institutional history with named actors, traceable funding, and measurable outcomes.
The doxonomic contribution is to treat it as a \emph{system} with formal structure, not merely as a sequence of political events.
\end{example}


\section{What Doxonomics Is Not}
\label{sec:ch1-is-not}

Several misreadings should be preempted.

\textbf{Doxonomics is not a claim that all beliefs are manufactured.}
Many beliefs arise from direct experience, logical inference, or scientific investigation.
Doxonomics studies the \emph{differential amplification} of beliefs: the fact that some beliefs receive vastly more institutional support than others.
It does not claim that evidence is irrelevant.

\textbf{Doxonomics is not a claim that truth is reducible to power.}
A belief can be both true and heavily funded, or false and poorly funded, or any combination.
The doxonomic audit reveals the material scaffolding of belief; the epistemic evaluation remains a separate task.
Chapter~\ref{ch:ontology-belief} develops this distinction in detail.

\textbf{Doxonomics is not a claim that people are passive.}
Audiences interpret, resist, reinterpret, ignore, and subvert the narratives they encounter.
A doxonomic system that spends billions may still fail to achieve capture.
Understanding why it fails is as important as understanding why it succeeds.
The field studies production, not automatic absorption.

\textbf{Doxonomics is not conspiracy theory.}
Conspiracy theories posit secret coordination by a small group.
Doxonomic systems can operate through open, legal, well-documented institutional processes.
The funding is often on public record (IRS 990 filings, lobbying disclosures, advertising expenditure reports).
The coordination is often structural (shared incentives, overlapping networks, common training) rather than secret.
The analogy is market competition: firms in the same industry produce similar products not because they conspire, but because they face similar incentives.


\section{Scope and Limits}
\label{sec:ch1-scope}

\begin{principle}[Scope of Doxonomics]
\label{prin:scope}
Doxonomics studies the \emph{differential amplification} of beliefs: the mechanisms by which some beliefs receive vastly more institutional support, repetition, and prestige than others, and the social consequences of this asymmetry.
\end{principle}

The scope is deliberately limited.
Doxonomics does not attempt to explain all belief formation.
It does not replace psychology (which studies individual cognition), epistemology (which studies the conditions for justified belief), or political theory (which studies the legitimacy of institutional arrangements).
It adds something that none of these fields provides: a systematic accounting of the material infrastructure of public belief.

A society that can estimate inflation, unemployment, and carbon intensity, but cannot estimate the industrial production of its own common sense, remains partially blind to itself.
Doxonomics seeks to reduce that blindness.


\section{Neighboring Fields and the Distinctive Contribution}
\label{sec:ch1-neighbors}

Doxonomics draws on and extends several existing disciplines, but is identical to none of them.

\begin{center}
\begin{tikzpicture}[
  field/.style={draw, rounded corners, minimum width=3.5cm, minimum height=0.7cm, align=center, font=\small},
  note/.style={font=\small, text width=6cm, align=left},
]
\matrix[row sep=0.5cm, column sep=0.4cm] {
  \node[field] {Political Economy}; &
  \node[note] {Studies material power but not its conversion into common sense.}; \\
  \node[field] {Sociology of Knowledge}; &
  \node[note] {Studies social formation of ideas but rarely performs financial accounting.}; \\
  \node[field] {Media Studies}; &
  \node[note] {Analyzes representation but not which representations are subsidized.}; \\
  \node[field] {Propaganda Studies}; &
  \node[note] {Maps persuasion techniques but often lacks formal models.}; \\
  \node[field] {Social Epistemology}; &
  \node[note] {Studies collective knowledge but not its material infrastructure.}; \\
  \node[field] {Behavioral Economics}; &
  \node[note] {Studies cognitive bias but not the institutional environment that exploits it.}; \\
  \node[field, fill=doxoblue!15] {\textbf{Doxonomics}}; &
  \node[note] {\textbf{Integrates all of the above around explicit accounting of material support, institutional routing, and measurable belief effects.}}; \\
};
\end{tikzpicture}
\end{center}

The distinctive contribution is the insistence on \emph{accounting}: tracing specific dollar amounts through specific institutions to specific narrative outputs with measurable effects on specific populations.
This is what makes doxonomics a potential science rather than a critical commentary.


\section{The Structure of This Book}
\label{sec:ch1-structure}

The book proceeds in five parts.

\textbf{Part~I (Foundations)} establishes the conceptual base: the ontology of belief, the intellectual prehistory, the core units of analysis, and the belief production chain.

\textbf{Part~II (Concepts and Models)} develops the formal apparatus: epistemic capital, narrative infrastructure, common-sense capture, doxonomic power, the economics of belief production, dynamical models, and sheaf-theoretic coherence.

\textbf{Part~III (Methods)} introduces the empirical toolkit: doxometric measurement, belief-flow accounting, network analysis, causal inference, text analysis, historical methods, experimental design, and model validation.

\textbf{Part~IV (Applications)} applies the machinery to specific domains: selfish-human-nature narratives, meritocracy, consumerism, inequality, race and coloniality, austerity, work discipline, national myth, technology, and crisis.

\textbf{Part~V (Resistance, Repair, and Reconstruction)} turns to counter-strategies, epistemic democracy, belief ecology design, ethics, and open problems.

Throughout, interactive \emph{playgrounds} allow the reader to manipulate funding parameters, simulate diffusion, and test model sensitivity.
These are available at the Piatra Institute's companion site.


\section{Notes and References}
\label{sec:ch1-notes}

The concept of belief as a social product has deep roots.
\textcite{marx1846} argued that the ruling ideas of any epoch are the ideas of the ruling class.
\textcite{gramsci1971} developed this into a theory of hegemony: the production of consent through civil society rather than through coercion alone.
\textcite{bourdieu1984} introduced the concept of \textit{doxa} as the unquestioned background of social life, extended in \textcite{bourdieu1991} to the analysis of language and symbolic power.
\textcite{berger1966} provided the foundational account of reality as a social construction.
\textcite{mannheim1936} developed the sociology of knowledge as a systematic discipline.
\textcite{althusser1970} formalized the concept of ideological state apparatuses.

Turning to media specifically, \textcite{lippmann1922} pioneered the analysis of how mass media manufacture public opinion.
\textcite{herman1988} provided the most systematic media-level account with the propaganda model.
\textcite{ellul1965} analyzed propaganda as a total social phenomenon, not merely a state tool.
The conversion problem connects to \textcite{bernays1928} on the engineering of consent.
For the economic logic of belief subsidies, see \textcite{coase1974} and \textcite{stigler1961}.

On the institutional history of specific belief production systems, \textcite{mirowski2009} traces the Mont P\`{e}lerin Society and the construction of neoliberal common sense.
\textcite{mayer2016} documents the dark-money networks behind American conservative ideology.
\textcite{oreskes2010} details the manufacture of doubt about climate science and tobacco harms.
\textcite{slobodian2018} provides a complementary genealogy of neoliberal internationalism.


\section{Exercises}

\begin{exercise}
\label{ex:ch1-identify}
Identify three widely held beliefs in your society that present themselves as descriptions of human nature or social reality (e.g., ``competition drives excellence,'' ``there is no alternative to the market'').
For each, list two institutions that regularly repeat them and estimate (order of magnitude) the annual budget of those institutions.
\end{exercise}

\begin{exercise}
\label{ex:ch1-equation}
Using the first-order belief equation~\eqref{eq:first-order}, argue qualitatively why a belief can persist even when the underlying evidence shifts against it, provided $\funding{k}$ and $\persistence{k}$ remain high.
Give a real-world example.
\end{exercise}

\begin{exercise}
\label{ex:ch1-chain}
Consider the belief ``markets are self-correcting.''
Map it onto the belief production chain, identifying at least one concrete actor or institution at each stage.
Where is the chain strongest?
Where is it most vulnerable to disruption?
\end{exercise}

\begin{exercise}
\label{ex:ch1-tuple}
Extend Definition~\ref{def:doxonomic-system} to include a \emph{resistance function} $\Omega: \beliefspace \to \R^m_{\geq 0}$ that models how population beliefs generate counter-pressure on funding channels.
Write down a fixed-point condition for equilibrium and discuss what it would mean for a society to be at such a fixed point.
\end{exercise}

\begin{exercise}
\label{ex:ch1-comparison}
Choose a belief that has \emph{lost} institutional support over the past century (e.g., the divine right of kings, explicit racial hierarchy as scientific consensus, the gold standard as natural law).
Trace its decline through the belief production chain.
At which stage did the chain break first?
\end{exercise}

\begin{exercise}
\label{ex:ch1-nonmonetary}
Identify a doxonomic system that operates primarily through non-monetary channels, such as family socialization, religious ritual, or architectural design.
How would you modify Definition~\ref{def:doxonomic-system} to accommodate non-monetary resource flows?
What does this reveal about the limits of a funding-centered framework?
\end{exercise}

\begin{exercise}
\label{ex:ch1-epidemiology}
The text compares doxonomics to epidemiology.
Push the analogy further: what would the doxonomic equivalents of ``pathogen,'' ``transmission vector,'' ``incubation period,'' ``herd immunity,'' and ``vaccination'' be?
Where does the analogy break down?
\end{exercise}
